\SourceLeaf{373}%
and finally they also have the same osculating circle if, in addition,
\[
\frac{rqq-2spq+tpp}{r'q'q'-2s'p'q'+t'p'p'}=\frac{p^3}{p'^3}=\frac{q^3}{q'^3}.
\]
This makes it easy to determine the osculating circle if one of the curves is assumed to be a circle.

Thus if $x,y$ are orthogonal coordinates, the center of the osculating circle is determined by the equations
\[
\begin{aligned}
x&=\xi+\frac{p(pp+qq)}{rqq-2spq+tpp}\\[5pt]
y&=\eta-\frac{q(pp+qq)}{rqq-2spq+tpp}.
\end{aligned}
\]
\needspace{7\baselineskip}
\begin{center}REMARK.\end{center}

The foregoing note is on the cover of \textit{Elementorum analyseos infinitorum Joanni Andreae Segneri Pars I.} Halle 1761 and most probably dates from the end of the eighteenth century.\par\nobreak\hfill\textsc{Stäckel.}\par
